\paper{数学类}

\begin{problems}

\item[一、]（本题 15 分）在空间直角坐标系中，设马鞍面 $S$ 的方程为 $x^{2}-y^{2}=2z$．设 $\sigma$ 为平面 $z=\alpha x+\beta y+\gamma$，其中 $\alpha,\beta,\gamma$ 为给定常数．求马鞍面 $S$ 上点 $P$ 的坐标，使得过 $P$ 且落在马鞍面 $S$ 上的直线均平行于平面 $\sigma$．

\item[二、]（本题 15 分）$A=\bigl(a_{ij}\bigr)_{n\times n}$ 为 $n$ 阶实方阵，满足
\begin{enumerate}
  \item $a_{11}=a_{22}=\cdots=a_{nn}=a>0$；
  \item 对每个 $i\ (i=1,2,\cdots,n)$，有 $\displaystyle\sum_{j=1}^{n}\abs{a_{ij}}+\sum_{j=1}^{n}\abs{a_{ji}}<4a$．
\end{enumerate}
求 $f(x_{1},\cdots,x_{n})=(x_{1},\cdots,x_{n})A\begin{pmatrix}x_{1}\\ \vdots\\ x_{n}\end{pmatrix}$ 的规范形．

\item[三、]（本题 20 分）元素皆为整数的矩阵称为整矩阵．设 $n$ 阶方阵 $A,B$ 皆为整矩阵．
\begin{enumerate}
  \item 证明以下两条等价：i) $A$ 可逆且 $A^{-1}$ 仍为整矩阵；ii) $A$ 的行列式的绝对值为 $1$；
  \item 若又知 $A,A-2B,A-4B,\cdots,A-2nB,A-2(n+1)B$，$A-2(n+n)B$ 皆可逆，且它们的逆矩阵皆为整矩阵．证明：$A+B$ 可逆．
\end{enumerate}

\item[四、]（本题 15 分）设 $f(x)$ 在 $[0,1]$ 上连续可微，在 $x=0$ 处有任意阶导数，
\[
  f^{(n)}(0)=0\quad(\forall n\ge 0),
\]
且存在常数 $C>0$ 使得
\[
  \abs{xf'(x)}\le C\abs{f(x)},\quad \forall x\in[0,1].
\]
证明：(1) $\lim\limits_{x\to 0^{+}}\dfrac{f(x)}{x^{n}}=0\ (\forall n\ge 0)$；(2) 在 $[0,1]$ 上成立 $f(x)\equiv 0$．

\item[五、]（本题 15 分）设 $\{a_{n}\},\{b_{n}\}$ 是两个数列，$a_{n}>0\ (n\ge 0)$，$\sum\limits_{n=1}^{\infty}b_{n}$ 绝对收敛，且
\[
  \frac{a_{n}}{a_{n+1}}\le 1+\frac{1}{n}+\frac{1}{n\ln n}+b_{n},\quad n\ge 2.
\]
求证：(1) $\dfrac{a_{n}}{a_{n+1}}<\dfrac{n+1}{n}\cdot\dfrac{\ln(n+1)}{\ln n}+b_{n}\ (n\ge 2)$；(2) $\sum\limits_{n=1}^{\infty}a_{n}$ 发散．

\item[六、]（本题 20 分）设 $f:\R\to(0,+\infty)$ 是一可微函数，且对所有 $x,y\in\R$，有
\[
  \abs{f'(x)-f'(y)}\le \abs{x-y}^{\alpha},
\]
其中 $\alpha\in(0,1)$ 是常数．

求证：对所有 $x\in\R$，有
\[
  \abs{f'(x)}^{\frac{\alpha+1}{\alpha}}<\frac{\alpha+1}{\alpha}f(x).
\]

\end{problems}
